\section{Write direct technical prose}
\label{sec:sentences}

Good technical prose lets the argument remain visible at sentence scale. Put
the main subject and verb near the front. The example papers use first-person
verbs to mark what the authors did. ``We compare'' names an experimental
choice, while ``we find'' reports its result. They reserve ``we hypothesize''
for an interpretation. Use the technical object as the subject when the
behavior belongs to the object.

\begin{quote}
\textbf{Author action:} We compare the two methods under a shared
initialization.

\textbf{Technical behavior:} Action chunking limits compounding error under
the stated stability condition.
\end{quote}

Direct prose does not require short sentences. A long sentence works when
every clause serves one relation. A short sentence fails when it adds drama
but no claim.

\subsection{Put conditions where they matter}

A fronted clause should change how the main claim is understood:

\begin{quote}
Although the estimator is biased, its error does not grow with the horizon in
this regime.

When the system is open-loop stable, action chunking gives the stated bound.

Because the critic supplies the terminal reward, the method can reuse
environment data off policy.
\end{quote}

Openings such as ``while'' and ``although'' become clutter when they provide
generic scene-setting. If the clause does not change how the main statement
should be read, lead with the statement.

\subsection{Make one main inference per sentence}

A sentence may contain a result and its essential condition or qualification.
Split it as soon as a second inference enters. The method, result, and proposed
explanation rarely belong in the same sentence.

\begin{quote}
OGPO reuses environment data through an off-policy critic. It spends additional
computation on denoising trajectories instead of collecting new environment
interactions.
\end{quote}

The two sentences describe related parts of the same idea, but each has one
main job. This separation also makes technical review easier. A coauthor can
challenge the data-reuse claim without untangling three other claims from the
same sentence.

\subsection{Give each paragraph one inference}

Most technical paragraphs make three or four moves:

\begin{enumerate}[leftmargin=*]
    \item state the claim;
    \item give the reason or technical distinction;
    \item present the theorem, comparison, measurement, or example;
    \item state the supported consequence or boundary.
\end{enumerate}

Use only the moves the paragraph needs. A two-sentence paragraph can be
complete. Vary sentence length and syntax so that paragraphs do not read as
filled-in forms. End on a consequence, limit, or transition that creates the
next question. A final sentence that merely repeats the first can usually be
deleted.

Paragraphs should also have explicit agents. Write ``we compare'' instead of
``a comparison is conducted,'' and ``the policy fails'' instead of ``a failure
is exhibited by the policy.'' Nominalizations hide the action and lengthen the
sentence. Direct verbs leave more room for the scientific detail.

\subsection{Report numbers in a fixed logical order}

For a quantitative result, lead with the comparator and metric. Then report the
value under the exact setting, including uncertainty, before interpreting it:

\begin{quote}
Compared with [baseline], [method] improves [metric] from [A] to [B] on
[setting], averaged over [seeds or trials]. This supports [narrow inference].
\end{quote}

Use the available numbers. ``Substantially better'' is less useful than the
measured difference. Reserve ``significant'' for a defined statistical claim.
If the paper reports uncertainty, say what it represents. If two methods use
different budgets, data, or initializations, expose that difference before the
reader reaches the result.

\subsection{Separate formal results from interpretation}

State the regime before a theorem's consequence:

\begin{quote}
Under [assumptions], [intervention] guarantees [bound].
\end{quote}

Then interpret the statement in a separate sentence:

\begin{quote}
Theorem~2 bounds [quantity] by [rate]. Thus, in [regime], the error does not
grow with [variable].
\end{quote}

Preserve every quantifier and condition. Keep the asymptotic rate and scope
exact. A clean sentence with a missing assumption is wrong. If an intuitive
explanation is not part of the theorem, label it as intuition rather than
folding it into the formal claim.

\subsection{Calibrate mechanism language}

Observation and explanation are different sentence types:

\begin{quote}
We observe [measured behavior]. We hypothesize that [mechanism] produces this
behavior because [support].
\end{quote}

Use ``shows'' or ``demonstrates'' for direct evidence. Use ``supports'' for
partial evidence. Use ``suggests'' or ``is consistent with'' for an
interpretation that remains uncertain. ``We believe'' does not supply evidence
for a hypothesis.

A controlled intervention can support a causal statement when the relevant
alternatives are held fixed. A visualization can reveal structure without
proving cause; the same is true of a correlation. The verb should make that
distinction visible.

\subsection{Use contrast to state scientific distinctions}

Contrast is useful when it identifies a real scientific distinction. Sample
efficiency differs from cheap computation; performance evidence differs from
mechanism evidence. State both sides in one clear relation.

Repeated formulas such as ``This is not X. It is Y'' or ``not only X but also
Y'' make the prose sound rehearsed. They also encourage false binaries. Write
the scientific distinction directly:

\begin{quote}
Current methods trade off sample efficiency and full-policy improvement.
\end{quote}

When evidence rules out one explanation but does not prove another, say both
facts without turning them into a slogan.

\subsection{Replace abstract wrappers with verbs}

Phrases such as ``the mechanism behind X,'' ``the role of Y in Z,'' and ``the
interplay between A and B'' name a relationship without saying what happens.
Use a direct verb when the evidence supports one:

\begin{quote}
\textbf{Abstract:} We study the mechanisms behind \emph{OGPO}'s stability.

\textbf{Direct:} We test whether conservative advantages prevent critic
exploitation.
\end{quote}

Sometimes the relationship itself is the question. State that question
precisely. ``We test whether stochastic supervision causes manifold
adherence'' is clearer than ``we study the role of stochasticity in manifold
adherence,'' and it does not claim that the test has already succeeded.

\subsection{Use lists only when the set matters}

A sentence is not difficult because it is long. It becomes difficult when
commas hold together several claims that have no hierarchy. The example papers
often use long sentences to develop one contrast or causal relation. Follow
that pattern.

Keep a full list when its members form a taxonomy, a set of assumptions, or a
checklist. In ordinary prose, decide how the items relate. Group them by role
or split the sentence:

\begin{quote}
\textbf{Packed:} We evaluate stability, robustness, efficiency, and
generalization.

\textbf{Grouped:} We first measure training stability and sample efficiency.
Separate deployment tests measure robustness and generalization.
\end{quote}

\subsection{Keep terms stable and use punctuation for grammar}

Introduce a technical term once, then give its operational meaning. Repeat the
canonical term instead of cycling through synonyms. Use a pronoun only when
its referent is obvious. Semantic LaTeX macros help keep method names and
symbols stable across authors and sections.

Use punctuation to expose the grammar. Commas mark real clause boundaries.
Semicolons join closely related claims; colons introduce definitions or lists.
Parentheses should contain secondary detail. Research questions belong in the
prose when the question organizes the paper, not as a routine rhetorical
setup. No sentence should depend on punctuation for drama.
